% GENERATED FILE. Edit canonical lessons, then run npm run generate:books.
% canonical-source-hash: fnv1a64:6b53d1c5471d487e
% canonical-lessons: MW-C49-kako, MW-C49-mamo, MW-C49-chhoro, MW-C49-chhori, MW-C49-admi

\chapter{Uncle, Uncle, Boy, Girl, Man}
\label{ch:mw-uncle-uncle-boy-girl-man}
\hlchaptermodality{49}

\begin{chapteropening}
\textbf{By the end of this chapter:} \emph{I can say an uncle, an uncle, a boy, a girl and a man.}

Complete the last lesson of chapter 49: I can say an uncle, an uncle, a boy, a girl and a man.
\end{chapteropening}

\section[\texorpdfstring{\mw{काको} (kāko)}{काको (kāko)}]{\mw{काको} (kāko) — an uncle (father's brother)}

\label{lesson:MW-C49-kako}



\begin{quote}
Before the new one: say the Marwadi for a son, then the Marwadi for a daughter.
\end{quote}

\subsection*{You'll want to know: \mw{काको}}
\textbf{\mw{काको}} — \emph{kāko} — \textquotedblleft{}an uncle (father's brother)\textquotedblright{}.

\begin{cousinweb}
\textbf{\mw{काको}} (\emph{kāko}) is a tadbhava, a word inherited and reshaped by sound change. It comes from \emph{*kākka}, \textquotedblleft{}senior male relative\textquotedblright{}, a form CDIAL reconstructs from its descendants. CDIAL gives no Marwari form in this entry (CDIAL 2998); the match is through Hindi \emph{kākā} and Gujarati \emph{kāko}. CDIAL takes the word from Dravidian: Kannada \emph{kakka}, \textquotedblleft{}uncle\textquotedblright{}.
\end{cousinweb}

\begin{grammarlens}[title={What you've built}]
One more word for this chapter.
\end{grammarlens}

\subsection*{Your turn}
\begin{itemize}
\raggedright
  \item \emph{Say it:} \emph{kāko}
  \item \emph{Say it:} \emph{kāko}, once more
  \item \emph{Say it:} say \emph{kāko}
  \item \emph{Recall:} say the Marwadi for a son, then the Marwadi for a daughter, then say \emph{kāko} again
\end{itemize}

\begin{tcolorbox}[breakable,colback=teal!4,colframe=teal!35!black,title={Before you move on}]
What does \mw{काको} mean? (\textquotedblleft{}An uncle (father's brother)\textquotedblright{}.) Say it once more.
\end{tcolorbox}
\section[\texorpdfstring{\mw{मामो} (māmo)}{मामो (māmo)}]{\mw{मामो} (māmo) — an uncle (mother's brother)}

\label{lesson:MW-C49-mamo}



\begin{quote}
Before the new one: say the Marwadi for a daughter, then the Marwadi for an uncle.
\end{quote}

\subsection*{You'll want to know: \mw{मामो}}
\textbf{\mw{मामो}} — \emph{māmo} — \textquotedblleft{}an uncle (mother's brother)\textquotedblright{}.

\begin{cousinweb}
\textbf{\mw{मामो}} (\emph{māmo}) is a tadbhava, a word inherited from Sanskrit and reshaped by sound change. It comes from \emph{māma}, \textquotedblleft{}uncle\textquotedblright{}. CDIAL gives no Marwari form in this entry (CDIAL 10055); the match is through Hindi \emph{māmā} and Gujarati \emph{māmo}. CDIAL calls it a nursery word and notes similar words in Dravidian.
\end{cousinweb}

\begin{grammarlens}[title={What you've built}]
One more word for this chapter.
\end{grammarlens}

\subsection*{Your turn}
\begin{itemize}
\raggedright
  \item \emph{Say it:} \emph{māmo}
  \item \emph{Say it:} \emph{māmo}, once more
  \item \emph{Say it:} say \emph{māmo}
  \item \emph{Recall:} say the Marwadi for a daughter, then the Marwadi for an uncle, then say \emph{māmo} again
\end{itemize}

\begin{tcolorbox}[breakable,colback=teal!4,colframe=teal!35!black,title={Before you move on}]
What does \mw{मामो} mean? (\textquotedblleft{}An uncle (mother's brother)\textquotedblright{}.) Say it once more.
\end{tcolorbox}
\section[\texorpdfstring{\mw{छोरो} (chhoro)}{छोरो (chhoro)}]{\mw{छोरो} (chhoro) — a boy}

\label{lesson:MW-C49-chhoro}



\begin{quote}
Before the new one: say the Marwadi for an uncle, then the Marwadi for an uncle.
\end{quote}

\subsection*{You'll want to know: \mw{छोरो}}
\textbf{\mw{छोरो}} — \emph{chhoro} — \textquotedblleft{}a boy\textquotedblright{}.

\begin{cousinweb}
\textbf{\mw{छोरो}} (\emph{chhoro}) is a tadbhava, a word inherited and reshaped by sound change. It comes from \emph{*chokara}, \textquotedblleft{}boy\textquotedblright{}, a form CDIAL reconstructs from its descendants. CDIAL gives no Marwari form in this entry (CDIAL 5070); the match is through Hindi \emph{chorā} and Gujarati \emph{choro}. CDIAL compares \emph{*chōka}, \textquotedblleft{}boy\textquotedblright{}.
\end{cousinweb}

\begin{grammarlens}[title={What you've built}]
One more word for this chapter.
\end{grammarlens}

\subsection*{Your turn}
\begin{itemize}
\raggedright
  \item \emph{Say it:} \emph{chhoro}
  \item \emph{Say it:} \emph{chhoro}, once more
  \item \emph{Say it:} say \emph{chhoro}
  \item \emph{Recall:} say the Marwadi for an uncle, then the Marwadi for an uncle, then say \emph{chhoro} again
\end{itemize}

\begin{tcolorbox}[breakable,colback=teal!4,colframe=teal!35!black,title={Before you move on}]
What does \mw{छोरो} mean? (\textquotedblleft{}A boy\textquotedblright{}.) Say it once more.
\end{tcolorbox}
\section[\texorpdfstring{\mw{छोरी} (chhorī)}{छोरी (chhorī)}]{\mw{छोरी} (chhorī) — a girl}

\label{lesson:MW-C49-chhori}



\begin{quote}
Before the new one: say the Marwadi for an uncle, then the Marwadi for a boy.
\end{quote}

\subsection*{You'll want to know: \mw{छोरी}}
\textbf{\mw{छोरी}} — \emph{chhorī} — \textquotedblleft{}a girl\textquotedblright{}.

\begin{cousinweb}
\textbf{\mw{छोरी}} (\emph{chhorī}) is a tadbhava, a word inherited and reshaped by sound change. It comes from \emph{*chokara}, \textquotedblleft{}boy\textquotedblright{}, a form CDIAL reconstructs from its descendants. CDIAL gives no Marwari form in this entry (CDIAL 5070); the match is through Hindi \emph{chorī} and Gujarati \emph{chorī}. It is the feminine of \emph{chhoro}, \textquotedblleft{}boy\textquotedblright{}.
\end{cousinweb}

\begin{grammarlens}[title={What you've built}]
One more word for this chapter.
\end{grammarlens}

\subsection*{Your turn}
\begin{itemize}
\raggedright
  \item \emph{Say it:} \emph{chhorī}
  \item \emph{Say it:} \emph{chhorī}, once more
  \item \emph{Say it:} say \emph{chhorī}
  \item \emph{Recall:} say the Marwadi for an uncle, then the Marwadi for a boy, then say \emph{chhorī} again
\end{itemize}

\begin{tcolorbox}[breakable,colback=teal!4,colframe=teal!35!black,title={Before you move on}]
What does \mw{छोरी} mean? (\textquotedblleft{}A girl\textquotedblright{}.) Say it once more.
\end{tcolorbox}
\section[\texorpdfstring{\mw{आदमी} (ādmī)}{आदमी (ādmī)}]{\mw{आदमी} (ādmī) — a man}

\label{lesson:MW-C49-admi}



\begin{quote}
Before the new one: say the Marwadi for a boy, then the Marwadi for a girl.
\end{quote}

\subsection*{You'll want to know: \mw{आदमी}}
\textbf{\mw{आदमी}} — \emph{ādmī} — \textquotedblleft{}a man\textquotedblright{}.

\begin{cousinweb}
\textbf{\mw{आदमी}} (\emph{ādmī}) is a loan from Arabic \emph{ādamī}, \textquotedblleft{}human, of Adam\textquotedblright{}, which reached Marwadi through Persian. Its Arabic root is \emph{ʾ-d-m} (Lane). It is the name \emph{Ādam} with the ending \emph{-ī}, \textquotedblleft{}belonging to\textquotedblright{}.
\end{cousinweb}

\begin{grammarlens}[title={What you've built}]
One more word for this chapter.
\end{grammarlens}

\subsection*{Your turn}
\begin{itemize}
\raggedright
  \item \emph{Say it:} \emph{ādmī}
  \item \emph{Say it:} \emph{ādmī}, once more
  \item \emph{Say it:} say \emph{ādmī}
  \item \emph{Recall:} say the Marwadi for a boy, then the Marwadi for a girl, then say \emph{ādmī} again
\end{itemize}

\begin{tcolorbox}[breakable,colback=teal!4,colframe=teal!35!black,title={Before you move on}]
What does \mw{आदमी} mean? (\textquotedblleft{}A man\textquotedblright{}.) Say it once more.
\end{tcolorbox}
